\section{Related Work}
\paragraph{Self-evolving memory and management.} ReasoningBank distills reusable reasoning from self-judged successful and failed experience~\citep{ouyang2026reasoningbank}; reward-conditioned construction and reward/reuse distortion are studied elsewhere~\citep{wang2025memalpha,asadolahi2026inflation,yang2026romerl}. \citet{xiong2026memorymanagement} analyze experience-following, error propagation, and quality-aware addition/deletion. C1 does not claim feedback-driven memory or generic governance as new; it creates two persistent states from the \emph{same} source trajectory and asks how far that controlled contrast survives.

\paragraph{Retrieval, utilization, and state-conditioned memory.} Mem2ActBench directly evaluates active memory application~\citep{shen2026mem2actbench}; MemCoRL co-optimizes retrieval and utilization~\citep{liu2026memcorl}; Chain-of-Memory strengthens post-retrieval utilization~\citep{xu2026chainofmemory}; and SAMem explicitly aligns fine-grained memory with the agent's current state~\citep{wang2026samem}. QCR and other diagnostics likewise separate retrieval from later reuse~\citep{li2026qcr,yuan2026retrievalutilization}. Thus ``retrieval is not use,'' state-aware memory, and utilization optimization are not our novelty. Our residual is controlled transport identification: preserve the intervention identity across $W/E/U/O/F$, determine what stage-specific evidence can justify, and use known-style post-retrieval transformations only as prospective probes of whether the diagnosed surface is engineering-relevant. The repair probes therefore validate or bound the diagnosis rather than introduce a new utilization architecture.
